\section{A Fano-type model of the tensor fibration}
\label{fano:section}

The preceding effectivity arguments correct a divisor on a smooth base.
Here we instead contract the base until no correction is needed there.
The additional geometric conclusion is a Fano-type model of this
auxiliary base.  Its contracted divisors pull back to exceptional
divisors over the original variety, so a principal correction on the
model still produces a section on the original variety.

We give a Fano-type proof of the tensor conversion in
Corollary~\ref{conv:tensor-conversion}.  Its input is a smooth
connected projective $X$ with smoothly semipositive $L=-K_X$, and
nonzero sections of $(\Omega_X^1)^{\otimes p}\otimes mL$ for one
fixed $p\geq0$ and unbounded positive $m$.  The distinctive output
is the supported birational model of Theorem~\ref{fano:model}.
We first construct the auxiliary base and verify the hypotheses of
the toroidal volume theorem.  We then identify exactly which base
divisors its log terminal model can contract.  The final passage to
an anticanonical section uses the principal divisors themselves,
not only their numerical classes.

\subsection{A fibration tangent to the null directions}

The cases $\dim X=0$ and $p=0$ are immediate.  Suppose otherwise.
Write $\mathcal E=(\Omega_X^1)^{\otimes p}$ and apply the
determinant extraction in Lemma~\ref{conv:determinant}.
Retain its determinant line $M$, rather than its negative: there is a nonzero map
\[
 M\longrightarrow\bigwedge^s\mathcal E
\]
for some $s>0$, and effective integral divisors
\begin{equation}\label{fano:determinants}
 N_m\sim M+mL
\end{equation}
for an unbounded set of positive integers $m$.  The generic wedges
used in that lemma are regular because they take values in the
saturated subsheaf and its reflexive determinant.  This is the
determinant method of \cite[Lemma~4.1]{LP2018} and
\cite[Lemma~5.1]{LMPTX2023}, with the unbounded exponents supplied
here as part of the hypothesis.

Write $n=\dim X$, let $\alpha$ be the normalized curvature of $L$,
and let $\nu$ be its maximum rank.  If $\nu=n$, then $L$ is nef
and $L^n=\int_X\alpha^n>0$, so $L$ is big and conversion follows.
Fix an ample class $A$ and assume $\nu<n$.  We will use
\begin{equation}\label{fano:null-degree}
 N_m\cdot[\alpha]^\nu A^{n-1-\nu}=0.
\end{equation}
Here is a differential-geometric proof of this equality.  For each
$\epsilon>0$, Yau's prescribed-Ricci theorem~\cite{Yau1978} gives a
K\"ahler metric in $[\alpha]+\epsilon A$ whose Ricci form equals
the prescribed semipositive form $\alpha$, up to the fixed positive
normalization.  K\"ahler symmetry
identifies the mean curvature of the tangent bundle with its
nonnegative Ricci endomorphism.  Thus the mean curvature of every
cotangent tensor power, and of its exterior powers, is nonpositive.
Consequently
\[
 c_1(M)([\alpha]+\epsilon A)^{n-1}\leq0.
\]
Indeed, a positive degree would allow a metric on $M$ with constant
positive mean curvature, by the scalar Poisson equation.  The induced
mean curvature on
$\bigwedge^s\mathcal E\otimes M^{-1}$ would then be
negative definite, and the integrated Bochner formula would forbid
its nonzero section.  Since $\alpha^{\nu+1}=0$, the coefficient of
the lowest possible power $\epsilon^{n-1-\nu}$ is a positive
constant times $c_1(M)[\alpha]^\nu A^{n-1-\nu}$ and is nonpositive.
The same intersection for $N_m$ equals it by
\eqref{fano:determinants}, and is nonnegative by effectivity.  This
proves \eqref{fano:null-degree}.

Choose a finite sum $D$ of the $N_m$, allowing repetitions, whose
Iitaka dimension $d$ is maximal among such sums.  These dimensions
are integers in $\{0,\ldots,n\}$, so a maximum exists.  Choose
$q>0$ for which $|qD|$ has image dimension $d$.  On the open set
where $\alpha$ has rank $\nu$, the differential of this rational
map annihilates $\ker\alpha$.  To see this, each member of $|qD|$
has zero intersection with $[\alpha]^\nu A^{n-1-\nu}$.  A smooth
local component whose tangent hyperplane failed to contain
$\ker\alpha$ would retain rank $\nu$ for the restricted form,
and hence contribute positive mass to that intersection.  Away from
the base locus, a projective hyperplane detecting a nonzero
differential on $\ker\alpha$ would produce exactly such a smooth
local component.  Hyperplanes detect the differential of the map,
so the assertion follows.  In particular $d\leq\nu<n$.

Resolve the base ideal by $p_1:X_1\to X$, with $X_1$ smooth and the
resolution an isomorphism off that ideal, and take the Stein
factorization $f_1:X_1\to B$.  The normal projective base $B$ has
relatively algebraically closed function field in $\mathbb C(X)$.
There are an ample divisor $A_B$ and an effective fixed divisor $I$
with
\begin{equation}\label{fano:fixed-free}
 p_1^*(qD)\sim f_1^*A_B+I.
\end{equation}
Every $p_1$-exceptional prime lies in $\operatorname{Supp}I$, since
its center lies in the base ideal and that ideal has positive order
on the prime.

Fix $m_1<m_2$ among the determinant exponents and put
\[
 G=\frac{N_{m_2}-N_{m_1}}{m_2-m_1}\sim_{\mathbb Q}L,
 \qquad J=N_{m_1}-m_1G.
\]
On any compatible higher model $\mu:W\to X$, $f:W\to Y$, with
$Y$ birational over $B$, there are $\mathbb Q$-principal divisors
$T_m$ on $Y$ such that
\begin{equation}\label{fano:interpolation}
 \mu^*N_m=m\mu^*G+\mu^*J+f^*T_m\qquad(m>m_2).
\end{equation}
For the principal difference of
\[
 (m_2-m_1)N_m+(m-m_2)N_{m_1}
 \quad\hbox{and}\quad(m-m_1)N_{m_2}
\]
is a rational function on $B$.  Otherwise it is transcendental over
$\mathbb C(B)$, by relative algebraic closedness, and adjoining it
to the map for $|qD|$ increases the image dimension.  The joint map
is given by products of sections in the system for the permitted
finite sum $qD+(m-m_1)N_{m_2}$, contradicting maximality.
Dividing its principal divisor by $m_2-m_1$ proves
\eqref{fano:interpolation}.  If $d=0$, these functions are constant:
$N_m=mG+J$, so every coefficient of $G$ is nonnegative by letting
$m\to\infty$.  This already proves conversion.  We henceforth
have $0<d<n$.

We next establish the adjoint rank required for the volume metric.
For each integer $\ell>0$, the general fiber of $f_1$ satisfies
\begin{equation}\label{fano:fixed-rank}
 h^0(X_{1,b},\mathcal O_{X_1}(\ell I)|_{X_{1,b}})=1.
\end{equation}
If the dimension were larger, generic base change and a sufficiently
ample base twist would give sections of $\ell I+t f_1^*A_B$
whose ratio varies on the general fiber.  For $t\gg\ell$, multiply
them by the section of $(t-\ell)I$.  They become sections of
$t p_1^*(qD)$.  The same complete system recovers the base using the
sections from $|tA_B|$ multiplied by $s_I^t$, so its image dimension
exceeds $d$, a contradiction.  The Jacobian divisor $K_{X_1/X}$
is effective and bounded by some $\ell I$.  Its nonzero canonical
section and \eqref{fano:fixed-rank} therefore give rank one for its
full direct image.  Further smooth modifications do not change this
rank: their extra relative canonical divisors are effective and
exceptional, and their section pushforwards equal $\mathcal O$ by
normality.  Thus on any smooth resolution $\pi:Z\to X$,
$g:Z\to Y$ dominating these models,
\begin{equation}\label{fano:adjoint-rank}
 \operatorname{rank}g_*\mathcal O_Z(K_{Z/X})=1.
\end{equation}
The Jacobian section identifies
$g_*(K_{Z/Y}+\pi^*L)$ with $-K_Y$ on the smooth base-change locus.
This is an ordinary rank assertion, involving the entire adjoint
space.

\subsection{The integrable correction and the model it permits}

Apply birational weak toroidalization~\cite[Theorem~1.1]{ADK2013},
marking the nonisomorphism locus over $X$, and equidimensional
toroidal subdivision~\cite[Proposition~4.4]{AK2000}.  We obtain
$\mu:W\to X$, $f:W\to Y$ with $W$ normal projective, $Y$ smooth
projective, and $f$ surjective, equidimensional and toroidal for
strict toroidal embeddings.  All modifications are birational; no
reduced-fiber alteration is required.  The source can be singular,
and its resolution is used only for integration.  The effective
exceptional Weil divisor $K_{W/X}$ is supported in the toroidal
boundary.

Choose a global K\"ahler form $\omega_Y$ on $Y$.  By generic
smoothness and cohomology-and-base-change, there is a nonempty
Zariski open $U\subset Y$ on which the resolved map $g$ is smooth
and its adjoint direct image is the line identified above.  The
kernel calculation persists off the exceptional loci, since the map
from the base to the original image is generically finite.  On a
small patch above $U$, in the maximum-rank region and away from
the Jacobian divisor, it gives
\[
 \ker\pi^*\alpha\subset\ker dg,
 \qquad \pi^*\alpha\geq c\,g^*\omega_Y\quad(c>0).
\]
The second inequality follows from the first by semipositive linear algebra
and a lower bound for the nonzero eigenvalues on a smaller patch.
Together with \eqref{fano:adjoint-rank}, these verify all hypotheses
of the toroidal volume theorem in
\cite[Theorem~2.1]{CompanionG}.

Explicitly, for a base prime $P$ and primes $S\subset W$ dominating
it, put $a_S=\operatorname{ord}_S K_{W/X}$ and
$e_S=\operatorname{ord}_S f^*P$.  The theorem supplies
\begin{equation}\label{fano:correction}
 C=\sum_P\max\left\{0,\min_{S\mapsto P}
           \left(\frac{1+a_S}{e_S}-1\right)\right\}P
\end{equation}
and a semipositive metric on $\Lambda=-K_Y+C$ whose local squared
frame norms are integrable and whose curvature is smooth and
strictly positive on a nonempty ordinary open set.  Moreover $C$
is effective rational with finite support and $f^*C$ is entirely
exceptional over $X$.  The companion theorem proves the all-valuation
integrability and the extensions across divisors and codimension
two; these conclusions use the precise ordinary rank and patch
verified here.

In particular $\kappa(Y,C)=0$.  Indeed, a rational function with
poles bounded by a positive divisible multiple of $C$ pulls back
to a function with only $\mu$-exceptional poles.  It extends across
$X$ by normality and is constant because $X$ is connected projective.
Also there is an effective big rational divisor
\begin{equation}\label{fano:klt-representative}
 \Delta\sim_{\mathbb Q}\Lambda,
 \qquad (Y,\Delta)\ \hbox{klt}.
\end{equation}
Apply Corollary~\ref{cor:big-klt-representative} to the unweighted
integrable metric on $\Lambda$.  Its proof first subtracts a small
ample rational class while keeping a klt representative, then adds
a general representative of that ample class.  Thus the boundary
here is effective, big and klt, exactly as required by the birational
theorem below.

\begin{theorem}[A model with supported contractions]\label{fano:model}
Let $Y$ be smooth projective.  Suppose $C\geq0$ is a rational divisor
with $\kappa(Y,C)=0$, and $\Delta\geq0$ is big and rational with
$(Y,\Delta)$ klt and $K_Y+\Delta\sim_{\mathbb Q}C$.  Then there is
a birational contraction $\chi:Y\dashrightarrow Y_0$ to a projective
$\mathbb Q$-factorial variety of Fano type.  Every prime contracted
by $\chi$ lies in $\operatorname{Supp}C$.  Every pseudoeffective
rational divisor on $Y_0$ is $\mathbb Q$-linearly equivalent to an
effective rational divisor.
\end{theorem}

\begin{proof}
The big-boundary log terminal model theorem
\cite[Theorem~1.1]{BCHM2010} applies: the pair is klt, its boundary
is effective and big, and its adjoint is pseudoeffective since it
is rationally equivalent to $C\geq0$.  It gives a projective
$\mathbb Q$-factorial log terminal model $\chi:Y\dashrightarrow Y_0$.
The map extracts no divisors.  With $\Delta_0=\chi_*\Delta$, the
transformed pair is klt and $K_{Y_0}+\Delta_0$ is nef.  On a common
resolution $a:U\to Y$, $b:U\to Y_0$, its negativity property
\cite[Definitions~3.6.1 and~3.6.6]{BCHM2010} gives
\[
 a^*(K_Y+\Delta)=b^*(K_{Y_0}+\Delta_0)+E,
\]
where $E\geq0$ is $b$-exceptional and contains the strict transform
of every $\chi$-contracted prime.  Choose a $\mathbb Q$-principal
divisor $T$ with $K_Y+\Delta=C+T$.  Using the same rational
functions and coefficients downstairs, and then pulling back,
cancels $T$ on both sides.  Therefore, with $C_0=\chi_*C$,
\begin{equation}\label{fano:negative-equation}
 a^*C=b^*C_0+E.
\end{equation}
Since $Y_0$ is $\mathbb Q$-factorial, $C_0$ is rationally Cartier;
it is effective, so its pullback is effective.  The strict transform
of each contracted prime has positive coefficient on the right
of \eqref{fano:negative-equation}, and hence on the left.  That
prime must belong to $\operatorname{Supp}C$.

The boundary $\Delta_0$ remains big under the birational contraction.
The nef adjoint is semiample by big-boundary base-point-freeness
\cite[Corollary~3.9.2]{BCHM2010}.  The effective exceptional
term $E$ preserves the spaces of sections of all sufficiently
divisible adjoint multiples, by normality.  Thus
\[
 \kappa(Y_0,K_{Y_0}+\Delta_0)=\kappa(Y,C)=0.
\]
A semiample rational divisor with Iitaka dimension zero is rationally
linearly trivial: a base-point-free multiple defines a map to a
point and is therefore trivial.  Hence
$K_{Y_0}+\Delta_0\sim_{\mathbb Q}0$.

Write $\Delta_0\sim_{\mathbb Q}A_0+N_0$ with $A_0$ ample rational
and $N_0\geq0$.  For a sufficiently small rational $\lambda>0$, the
boundary $\Theta=(1-\lambda)\Delta_0+\lambda N_0$ is effective and
klt; this follows by checking the finitely many coefficients on a
common log resolution.  We obtain
\[
 -(K_{Y_0}+\Theta)\sim_{\mathbb Q}\lambda A_0,
\]
which is ample.  This proves that $Y_0$ is of Fano type.

Finally, let $B_0$ be a pseudoeffective rational divisor on $Y_0$.
For a small rational $t>0$, the class
$tB_0-(K_{Y_0}+\Theta)$ is ample.  Choose a general effective
rational representative $U_0$ of that class with sufficiently small
coefficients so that $(Y_0,\Theta+U_0)$ remains klt.  Its boundary
is effective and big, and its adjoint is rationally equivalent to
$tB_0$, which is pseudoeffective.  Big-boundary nonvanishing
\cite[Theorem~D]{BCHM2010} makes this adjoint real-linearly effective.
Lemma~\ref{lem:rational-recovery} makes it rationally linearly
effective.  Division by $t$ proves the assertion for $B_0$.
\end{proof}

\subsection{Returning the correction to the original variety}

Apply Theorem~\ref{fano:model} to \eqref{fano:klt-representative}.
We now finish this proof of Corollary~\ref{conv:tensor-conversion}.  For each base
prime $P$ set
\[
 g_P=\min_{S\mapsto P}\frac{\operatorname{coeff}_S\mu^*G}{e_S},
 \qquad
 j_P=\max_{S\mapsto P}\frac{\operatorname{coeff}_S\mu^*J}{e_S},
 \qquad B_g=\sum_Pg_PP,\quad B_j=\sum_Pj_PP.
\]
These are finite rational sums.  The horizontal coefficients of
$\mu^*G$ are nonnegative by \eqref{fano:interpolation} and
$m\to\infty$.  At a component attaining the minimum $g_P$, that
same equation and effectivity of $\mu^*N_m$ give
\[
 m g_P+j_P+\operatorname{coeff}_P T_m\geq0.
\]
Thus $mB_g+B_j+T_m\geq0$.  Push this inequality by strict transform
to $Y_0$.  Principal terms remain principal because $\chi$ extracts
no divisors.  Dividing numerical classes by $m$ proves that
$(B_g)_0$ is pseudoeffective.  Theorem~\ref{fano:model} supplies
a $\mathbb Q$-principal adjustment making $(B_g)_0$ effective.
Use its same rational functions and coefficients on $Y$, and call
the resulting principal divisor $T_*$.  Then $B_g+T_*$ is effective
outside the primes contracted by $\chi$.

The divisor $\mu^*G+f^*T_*$ has nonnegative coefficient at every
prime of $W$ that is not exceptional over $X$.  Horizontals were
already treated.  Above an uncontracted prime this follows from
the defining minimum $g_P$.  Above a contracted prime, every
component is exceptional over $X$, since that prime lies in
$\operatorname{Supp}C$ and $f^*C$ is exceptional.  Equidimensionality
leaves no other vertical primes.  Pushing forward to $X$ therefore
gives an effective rational divisor.  Its class is
$\mathbb Q$-linearly equivalent to $G\sim_{\mathbb Q}L$, because
the base principal adjustment pulls back and pushes forward as a
principal divisor.  Clearing denominators completes the proof.

The Fano-type conclusion belongs to $Y_0$, the model of an auxiliary
linear-system base.  The argument on $X$ uses exactly the supported
exceptional correction; it does not assert that $X$ itself is of
Fano type.
